\scanpage{159}
Finalement, on trouve des formules plus jolies en éliminant les $t_i$ et $\varepsilon_i$, et en ne travaillant qu'avec les $\sigma_i$. Cela nous amène à un contexte différent : celui des \deleted{\unsure{var.}} modules \inserted{involutifs (i.e. modules} munis d'une involution). \unsure{Éventuellemt|Éventuellement}, on se donne dans $E_i$ un sous-module $E_i^0$, de telle façon que $\sigma_i$ soit $-\mathrm{id}_{E_i^0}$ sur $E_i^0$, et soit id sur $E_i/E_i^0$ (cet $E_i^0$ est \deleted{\unsure{g}} unique \deleted{\unsure{s}} si $2$ inv., alors c'est le sous-module propre $E_i^-$, formé des \unsure{à|$x$} tels que $\sigma_i x=-x$) Un \uline{accouplement involutif} \deleted{\unsure{de ces|de}} \inserted{vers} un module $E$, sera une application multilinéaire
\diagram{\begin{tikzcd}[ampersand replacement=\&,row sep=small]
\deleted{\unsure{F}}E_1\times\cdots\times E_n\arrow[r,"f"] \& E\quad\text{munie de}\\
\deleted{\unsure{f^0:f}}\ E_1^0\times\cdots\unsure{+|\times} E_n^0\arrow[ur,"f^0"'] \&
\end{tikzcd}}
satisfaisant les relations de la forme
\[
 \begin{aligned}
 f(x_1+\alpha_1,x_2\ldots \unsure{a_n|x_n})&=f(x_1,\ldots x_n)+f^0(\alpha_1,\sigma_2(x_2)-x_2,\\
 &\hspace{12em}\ldots\sigma_n(x_n)-x_n)
 \end{aligned}
\]
pour $\alpha_1\in E_1^0$ (et les $n-1$ autres relations similaires)

On aura donc
\[
 f(\alpha_1,\alpha_2\ldots\alpha_n)=(-2)^{n-1}f^0(\alpha_1,\alpha_2,\ldots\alpha_n)
\]
% archive sheet number: 158

